% Note for collaborators: the recoverable-redundancy formalism (pages 1-2) and
% the experimental program (page 3). ICLR 2027 template, no title block.
% Compile: pdflatex.
\documentclass{article}
\usepackage{iclr2027_conference,times}
\iclrfinalcopy
\usepackage{amsmath,amssymb}
\usepackage{enumitem}

\begin{document}

\paragraph{Setup.}
A deployed teacher policy $\pi_T$ repeatedly encounters near-identical
situations. We store its successful rollouts in a library
$\mathcal{L} = \{(k_i, a_i)\}$, where each key $k_i = \varphi(o_i)$ is an
intermediate representation extracted from the teacher's forward pass and each
payload $a_i$ is the executed action chunk. At each decision step a scheduler
$\sigma$ selects either \textsc{replay} (execute the nearest neighbor's
payload) or \textsc{infer} (run the full forward pass), inducing a mixed policy
$\pi_\sigma$. The episode outcome is $Y \in \{0,1\}$ and the teacher-only
success rate is $\mathrm{SR}_T$. The paper formalizes the question of how many
steps can be replayed without degrading success.

\paragraph{Concept.}
A step of inference is redundant when the action it produces is already
determined by stored experience, so that the forward pass contributes no
information beyond what the library holds. This is analogous to inter-frame
compression in video coding: a \textsc{miss} step corresponds to an I-frame,
computed in full, and a \textsc{hit} step to a P-frame, reconstructed from
stored data.

\paragraph{Step-level redundancy.}
Let $d_{\pi_T}$ denote the state distribution induced by the teacher acting
alone, and let $s_t \sim d_{\pi_T}$. Step $t$ is $\delta$-redundant if
\begin{equation}
\bigl|\,P\bigl(Y{=}1 \mid \mathrm{do}(a_t = \hat{a}_t^{\mathrm{replay}})\bigr)
      - P\bigl(Y{=}1 \mid \mathrm{do}(a_t = a_t^{\pi_T})\bigr)\,\bigr| \;\le\; \delta .
\label{eq:step}
\end{equation}
Since $Y$ is binary, the left-hand side equals the total variation distance
between the two outcome distributions. The definition is stated in terms of the
outcome rather than action similarity, because similar actions may lead to
failure and dissimilar actions may both succeed. The $\mathrm{do}(\cdot)$
operator denotes intervention; conditioning on the observed action would
introduce confounding, because the teacher's choice of action is itself
informative about the state. Eq.~\ref{eq:step} defines redundancy; the
operational quantity is Eq.~\ref{eq:R}.

\paragraph{Remark (failure of composition).}
Eq.~\ref{eq:step} is evaluated under $d_{\pi_T}$. Under deployment, earlier
steps may already have been replayed, so the state distribution at step $t$
differs from $d_{\pi_T}$ and replacement errors accumulate. Per-step redundancy
therefore does not compose: individual replaceability of each step does not
imply joint replaceability. The operational quantity must be defined at the
level of the scheduler and estimated from closed-loop rollouts. Offline metrics
computed on a fixed set of teacher trajectories cannot estimate it.

\paragraph{Recoverable redundancy.}
The replay share of a scheduler is
\begin{equation}
h(\sigma) \;=\; \mathbb{E}_{\pi_\sigma}\!\left[\frac{1}{T}\sum_{t=1}^{T}
\mathbb{1}\{\sigma_t = \mathrm{replay}\}\right],
\label{eq:h}
\end{equation}
where the expectation is taken under the trajectory distribution induced by
$\pi_\sigma$ itself, so that closed-loop effects are part of the definition.
For a class of schedulers $\mathcal{C}$, define
\begin{equation}
R_{\mathcal{C}}(\varepsilon) \;=\; \sup_{\sigma \in \mathcal{C}} \; h(\sigma)
\qquad \text{s.t.} \qquad
\mathrm{SR}(\pi_\sigma) \;\ge\; \mathrm{SR}_T - \varepsilon ,
\label{eq:R}
\end{equation}
where $\varepsilon$ is the admissible loss in success rate; the supremum need
not be attained. Since $R_{\mathcal{C}}$ is a supremum, the replay share of any
feasible scheduler is a lower bound on it, and reported values cannot overstate
the quantity.

\paragraph{Stratification by information set.}
Let $\mathcal{C}(\mathcal{I})$ denote the class of causal schedulers whose
decisions depend only on the information set $\mathcal{I}$, i.e.\ on present
and past quantities but not on future states or the outcome; without this
restriction the supremum is attained by schedulers that condition on $Y$ and
Eq.~\ref{eq:R} is uninformative. Consider
$\mathcal{I}_\varphi \subseteq \mathcal{I}_{\mathrm{sim}} \subseteq
\mathcal{I}_{\mathrm{all}}$, where $\mathcal{I}_\varphi$ contains the
quantities visible to $\varphi$-key retrieval, and
$\mathcal{I}_{\mathrm{sim}}$ additionally contains the simulator state. The
first inclusion holds because, in simulation, the state determines the
observation and hence the key. Then
\begin{equation}
h(\sigma_{\mathrm{thr}}) \;\le\; R_{\mathcal{C}(\mathcal{I}_\varphi)}(\varepsilon)
\;\le\; R_{\mathcal{C}(\mathcal{I}_{\mathrm{sim}})}(\varepsilon)
\;\le\; R_{\mathcal{C}(\mathcal{I}_{\mathrm{all}})}(\varepsilon) = R(\varepsilon).
\label{eq:sandwich}
\end{equation}
The first inequality holds because the threshold scheduler
$\sigma_{\mathrm{thr}}$ is an element of $\mathcal{C}(\mathcal{I}_\varphi)$;
the remaining inequalities hold because enlarging the information set enlarges
the feasible set. An empirical comparison between a measured oracle arm and a
measured threshold arm does not follow from Eq.~\ref{eq:sandwich}: the two
measurements are lower bounds on different terms.

\paragraph{Attribution.}
\begin{equation}
R \;=\; R\bigl(\varepsilon;\; \mathcal{M},\; \pi_T,\; \mathcal{L},\; \varphi\bigr)
\label{eq:quad}
\end{equation}
depends jointly on the task $\mathcal{M}$, the teacher $\pi_T$, the library
$\mathcal{L}$, and the key map $\varphi$. Reported values must specify all four
arguments; a statement of the form ``this task is $60\%$ redundant'' is not
well defined.

\paragraph{Monotonicity.}
For any scheduler class closed under ignoring library entries, for example
$\mathcal{C}(\mathcal{I}_{\mathrm{all}})$,
\begin{equation}
\mathcal{L} \subseteq \mathcal{L}' \;\Longrightarrow\;
R(\varepsilon;\mathcal{L}) \;\le\; R(\varepsilon;\mathcal{L}') .
\label{eq:mono}
\end{equation}
\emph{Proof.} A scheduler feasible on $\mathcal{L}$ remains feasible on
$\mathcal{L}'$ by ignoring the entries of
$\mathcal{L}' \setminus \mathcal{L}$, with identical success rate and replay
share; the feasible set therefore grows and the supremum does not decrease.
\hfill$\square$\\
The argument requires closure under ignoring entries and does not apply to
threshold rules: an added entry may attain a high retrieval score while being a
poor match, which changes the induced policy and can reduce the success rate.
Monotonicity under threshold rules is therefore an empirical question.

\paragraph{Intrinsic limit.}
\begin{equation}
R^{*}(\varepsilon) \;=\; \sup_{\mathcal{L}}\; R(\varepsilon; \mathcal{L})
\label{eq:intrinsic}
\end{equation}
depends only on the task and the teacher and upper-bounds
$R(\varepsilon;\mathcal{L})$ for every library. A step can be recognized as
redundant only if the library contains a sufficiently close experience, so
increasing $\mathcal{L}$ increases the recoverable share while
$R^{*}$ is unchanged. $R^{*}$ is not measurable, since the space of libraries
cannot be exhausted; it serves as a reference point and is not estimated.

\paragraph{Relation to calibration.}
Eq.~\ref{eq:R} also describes the system's calibration procedure. Selecting
the threshold on a calibration split to maximize the replay share subject to
the success-rate constraint solves Eq.~\ref{eq:R} empirically within
$\mathcal{C}(\mathcal{I}_\varphi)$; the threshold sweep of experiment E1
therefore traces the empirical Pareto frontier of
$R_{\mathcal{C}(\mathcal{I}_\varphi)}(\varepsilon)$. The deployment stages
correspond to the same program: collection constructs $\mathcal{L}$, shadow
operation estimates the score distributions, calibration solves the program,
and serving executes the resulting scheduler.

\newpage

\paragraph{Experimental program.}
Each experiment addresses one argument of Eq.~\ref{eq:quad} or one term of
Eq.~\ref{eq:sandwich}.
\begin{itemize}[leftmargin=1.3em, itemsep=0.35em, topsep=0.3em]
\item \textbf{E1, main frontier} (LIBERO complete; RoboCasa365 in progress).
Threshold sweeps measure $h(\sigma_{\mathrm{thr}})$ as a function of
$\varepsilon$, the first term of Eq.~\ref{eq:sandwich}; two benchmarks and two
teachers vary the $\mathcal{M}$ and $\pi_T$ arguments of Eq.~\ref{eq:quad}.
Current results: $50\%$ of decisions replayable at $<$1\,pp loss in success
rate and $63\%$ at $<$5\,pp on LIBERO-Spatial; $30\%$ at $\sim$1\,pp on
LIBERO-10. The difference between the suites indicates that recoverable
redundancy depends on task difficulty, consistent with Eq.~\ref{eq:quad}.
\item \textbf{E2, oracle arm} (to run). The same system with retrieval by
simulator state estimates the second term of Eq.~\ref{eq:sandwich}. It
separates coverage (the library lacks a close experience, so no key can
recover the step) from key loss (the experience is present but the index
misses it). If the oracle and threshold arms are close, the internal keys lose
little information; a large gap quantifies the margin attributable to the
index; a low oracle value indicates that the remaining gap is due to coverage.
\item \textbf{E3, library-size scaling} (complete). $R$ increases with
$\mathcal{L}$, from $0.27$ to $0.52$ on LIBERO-10 as the library grows from 1
to approximately 39 trajectories per task. This is the empirical counterpart
of Eq.~\ref{eq:mono}, which does not imply it for threshold rules; the
observed monotone increase is a finding about robustness to library growth.
Two boundaries: pure replay remains below the teacher at the full data budget,
so replay alone is not a viable policy; and no plateau is resolved, so
saturation is not claimed.
\item \textbf{E4, temporal structure} (complete). Hits are not independent
across steps. With the stickiness coefficient
$s = P(\mathrm{hit}_t \mid \mathrm{hit}_{t-1}) - P(\mathrm{hit})$, which is
zero under independence, we measure
$P(\mathrm{hit} \mid \mathrm{hit}) = 0.91$--$0.96$ against marginal hit rates
of $0.28$--$0.70$ across seven configurations; median replay runs last
$9$--$12$ steps (P90 up to 52); and $93$--$98\%$ of consecutive hits within a
run retrieve the same library trajectory. The last property explains why long
replay runs are compatible with the failure of composition: the replayed
segment follows a trajectory that previously succeeded, which bounds the
deviation introduced by replacement.
\item \textbf{E5, score discriminability} (offline; to run). The AUROC of the
retrieval score with respect to episode outcome measures whether the quantity
the threshold operates on is informative. The result is descriptive only:
episode-level labels paired with step-level scores establish association, not
a counterfactual.
\item \textbf{E6, cross-scene transfer} (in progress). Varies $\mathcal{M}$
with $(\pi_T, \mathcal{L}, \varphi)$ held fixed: the library is built in one
kitchen scene and queried in a held-out scene. Preliminary observations show
robustness on large-motion furniture tasks and weakness on precise dynamic
tasks; the controlled experiment has not yet been run.
\item \textbf{E7, key ablations} (offline replay of logged data). Varies
$\varphi$: internal representations against external encoders, extraction
depth (vision tower against mid-backbone), and multi-modal fusion (per-field
weighting against a single concatenated key). This addresses whether the
teacher's own representation is an adequate index.
\item \textbf{E8, warm start} (stale data; to re-run). Extends the scheduler's
action space from the binary choice in Eq.~\ref{eq:h} to a third option: a
partially redundant step initializes the flow integration with the retrieved
action and skips the early denoising steps. The control compares the full
teacher, blind step reduction, and warm start at an equal number of skipped
steps.
\item \textbf{E10, trained-router comparison} (in progress). A learned router
is an element of $\mathcal{C}(\mathcal{I}_\varphi)$, so its replay share is a
further lower bound on the same term of Eq.~\ref{eq:sandwich}. The comparison
asks whether the threshold rule is already close to the reachable set of its
class; if the learned scheduler does not raise the bound appreciably, the
threshold is empirically near-optimal within
$\mathcal{C}(\mathcal{I}_\varphi)$.
\end{itemize}
All reported values are lower bounds by construction (Eq.~\ref{eq:R}); the
gaps in Eq.~\ref{eq:sandwich} quantify the margin available to better
schedulers.

\end{document}
